% Préambule commun aux manuels de Pipit (français et anglais).
% Compilation : pdflatex manuel-fr.tex (deux fois), idem manuel-en.tex, depuis docs/src/.
\documentclass[11pt,a4paper]{article}
\usepackage[utf8]{inputenc}
\usepackage[T1]{fontenc}
\usepackage{lmodern}
\usepackage{microtype}
\usepackage[a4paper,margin=2cm,headheight=24pt]{geometry}
\usepackage{graphicx}
\usepackage[table]{xcolor}
\usepackage[most]{tcolorbox}
\usepackage{enumitem}
\usepackage{booktabs}
\usepackage{array}
\usepackage{tabularx}
\usepackage[nobottomtitles*]{titlesec}
\usepackage{needspace}
\usepackage{fancyhdr}
\usepackage{fontawesome5}
\usepackage{etoolbox}
\DeclareUnicodeCharacter{2192}{\ensuremath{\rightarrow}}
\DeclareUnicodeCharacter{2191}{\ensuremath{\uparrow}}
\DeclareUnicodeCharacter{2193}{\ensuremath{\downarrow}}
\DeclareUnicodeCharacter{25B2}{\ensuremath{\blacktriangle}}
\DeclareUnicodeCharacter{25BC}{\ensuremath{\blacktriangledown}}
\DeclareUnicodeCharacter{21BA}{\ensuremath{\circlearrowleft}}
\DeclareUnicodeCharacter{21BB}{\ensuremath{\circlearrowright}}
\usepackage{amssymb}

% Palette de l'application : vert olive, rouge brique pour les alertes
\definecolor{bordeaux}{HTML}{3A5420}
\definecolor{bordeauxl}{HTML}{EEF4E4}
\definecolor{alert}{HTML}{B5382B}
\definecolor{alertl}{HTML}{F8E6E3}
\definecolor{or}{HTML}{B8860B}
\definecolor{orl}{HTML}{FBF3DC}
\definecolor{sapin}{HTML}{2C7A55}
\definecolor{sapinl}{HTML}{E2F1E9}
\definecolor{sky}{HTML}{4E6E2A}
\definecolor{skyl}{HTML}{EEF4E4}
\definecolor{summit}{HTML}{6B7A8F}
\definecolor{night}{HTML}{2A2420}
\definecolor{ink}{HTML}{1F2430}
\AtBeginDocument{\color{ink}}

\titleformat{\section}
  {\Large\bfseries\sffamily\color{bordeaux}}
  {\makebox[1.6em][l]{\thesection}}{0.3em}{}
  [{\vspace{1pt}\color{bordeaux!35}\titlerule[1.3pt]}]
\titlespacing*{\section}{0pt}{14pt}{7pt}
\titleformat{\subsection}{\large\bfseries\sffamily\color{night}}{\thesubsection}{0.5em}{}
\pretocmd{\section}{\needspace{10\baselineskip}}{}{}
\pretocmd{\subsection}{\needspace{7\baselineskip}}{}{}
\widowpenalty=10000
\clubpenalty=10000
\raggedbottom
\setlength{\parindent}{0pt}
\setlength{\parskip}{6pt}
\setlist[enumerate]{itemsep=3pt,leftmargin=6mm}
\setlist[itemize]{itemsep=2pt,leftmargin=5mm}

\pagestyle{fancy}
\fancyhf{}
\renewcommand{\headrulewidth}{0.4pt}
\renewcommand{\headrule}{\hbox to\headwidth{\color{bordeaux!40}\leaders\hrule height \headrulewidth\hfill}}
\fancyfoot[C]{\sffamily\small\color{summit}\thepage}
\renewcommand{\sectionmark}[1]{\markboth{#1}{}}

\newcommand{\menu}[1]{\textsf{\textbf{\textcolor{bordeaux}{#1}}}}
\newcommand{\touche}[1]{\tcbox[on line,boxsep=0pt,left=2pt,right=2pt,top=1.5pt,bottom=1.5pt,
  colback=white,colframe=summit,boxrule=0.5pt,arc=1.5pt]{\footnotesize\textsf{#1}}}
\newcommand{\fichier}[1]{\texttt{\small #1}}

\newcommand{\boite}[5]{\newtcolorbox{#1}{enhanced,unbreakable,colback=#2,colframe=#3,
  coltitle=white,fonttitle=\bfseries\sffamily\small,boxrule=0.9pt,arc=2.5pt,
  left=2.5mm,right=2.5mm,top=1.2mm,bottom=1.2mm,title={#4~~#5}}}
\newcommand{\capture}[2][1]{\par\needspace{8\baselineskip}\begin{center}\setlength{\fboxsep}{0pt}\setlength{\fboxrule}{0.6pt}%
  \fcolorbox{summit!50}{white}{\includegraphics[width=\dimexpr#1\linewidth-1.3pt\relax]{#2}}\end{center}}

\usepackage[english]{babel}
\usepackage[colorlinks=true,linkcolor=bordeaux,urlcolor=bordeaux]{hyperref}
\hypersetup{pdftitle={Pipit — user guide},pdfauthor={},pdfcreator={Pipit}}
\boite{astucebox}{orl}{or!85!black}{\faLightbulb}{Good to know}
\boite{securitebox}{sapinl}{sapin}{\faShield*}{What Pipit never does}
\boite{attentionbox}{alertl}{alert}{\faExclamationTriangle}{Warning}
\newcommand{\astuce}[1]{\par\needspace{5\baselineskip}\begin{astucebox}#1\end{astucebox}}
\newcommand{\attention}[1]{\par\needspace{5\baselineskip}\begin{attentionbox}#1\end{attentionbox}}
\fancyhead[L]{\sffamily\small\color{bordeaux}\textbf{PIPIT — USER GUIDE}}
\fancyhead[R]{\sffamily\small\color{summit}\leftmark}
\input{version}

\begin{document}

\begin{titlepage}
\centering
\vspace*{2.2cm}
\includegraphics[width=3.6cm]{../../logo_512.png}\par
\vspace{8mm}
{\sffamily\Huge\bfseries\color{bordeaux} PIPIT\par}
\vspace{3mm}
{\sffamily\LARGE\color{ink} Safely remove USB drives and memory cards\par}
\vspace{2mm}
{\sffamily\large\color{summit} User guide \textbullet{} version \versionpipit\par}
\vspace{12mm}
\begin{tcolorbox}[enhanced,width=0.82\linewidth,colback=skyl,colframe=sky,arc=4pt,boxrule=1.2pt,left=5mm,right=5mm,top=3mm,bottom=3mm]
\sffamily\centering
The pipit is a small, discreet meadow bird:\\
Pipit stays next to the clock and only shows up when called.\\[2mm]
\textbf{Click the icon} \quad → \quad \textbf{click the device} \quad → \quad \textbf{unplug}
\end{tcolorbox}
\vfill
{\footnotesize\color{summit}\sffamily Windows 10 and 11 \textbullet{} free software (MIT licence)\par}
\end{titlepage}

\tableofcontents
\newpage

\section{What Pipit is for}

Unplugging a USB stick or an external drive while Windows is still writing to it can
damage files, or even the whole drive. Windows does have its “Safely Remove Hardware”
icon, but it is hidden, does not say \emph{which} program is in the way when it
refuses, and has no keyboard shortcut.

Pipit does the same job, better organised:

\begin{itemize}
\item an icon next to the clock: one click shows the list, one click on a line ejects;
\item a keyboard shortcut that works everywhere (\touche{Win}+\touche{Alt} by default);
\item when Windows refuses, the name of the program in the way, and a button to close it;
\item neither the Windows disk nor internal disks are ever offered.
\end{itemize}

Pipit is free and will stay free: no licence key, no “Pro” edition, no paid upgrade.

\section{Installing}

\begin{enumerate}
\item Download \fichier{Pipit-Setup.exe} from
  \url{https://github.com/ARP273-ROSE/Pipit/releases/latest}.
\item Run it: it installs for your account, without an administrator prompt.
\item On first start, Pipit asks two questions: whether to send incident reports, and
  whether to start with Windows. Answer \menu{Yes} to the second: that is what keeps it
  available whenever you plug in a drive.
\end{enumerate}

\astuce{If the icon does not show next to the clock, click the small \textbf{\^{}} arrow
in the taskbar: Windows puts new icons there. Drag Pipit's icon onto the taskbar to keep
it at hand.}

\section{Removing a device}\label{sec:retirer}

Click Pipit's icon (the green “eject” symbol) or press \touche{Win}+\touche{Alt}. The
menu opens under the mouse:

\capture[0.62]{../capture_menu-en.png}

Each line shows the drive letter, its name and its size. \textbf{One click on the line
ejects the device.} When it is done, a “Safe to remove” balloon appears: you can pull
the cable.

\begin{itemize}
\item \menu{Eject all} removes everything in the list at once.
\item \menu{Open in File Explorer} opens a drive's contents.
\item \menu{Settings…} opens the window described in section~\ref{sec:reglages}.
\end{itemize}

\section{When Windows refuses}\label{sec:refus}

Windows refuses to eject while a program keeps a file open on the device: a Word
document, a photo in a viewer, a paused video… Pipit first tries again a second later
(an antivirus scanning the stick often lets go on its own), then, if it still fails,
says who is in the way:

\capture[0.6]{../capture_refus-en.png}

\begin{itemize}
\item \menu{Close …} politely asks the program to close, exactly like the cross of its
  window: if it has an unsaved document, it will offer to save it. Pipit then tries
  again on its own.
\item \menu{Try again} retries the ejection once you have closed what was needed.
\end{itemize}

Special cases:
\begin{itemize}
\item \textbf{File Explorer}: a window is open on the drive, or the preview pane shows one
  of its files. Close that window; Pipit never closes File Explorer, which would take the
  desktop and the taskbar with it;
\item \textbf{an antivirus or Windows indexing}: wait a few seconds and try again;
\item \textbf{no program named}: Windows did not record it. Close the documents opened
  from the device.
\end{itemize}

\begin{securitebox}
\begin{itemize}[leftmargin=4mm]
\item force an ejection: if Windows refuses, nothing is removed;
\item kill a program (your unsaved documents are protected);
\item close File Explorer, the antivirus or a Windows program;
\item offer the Windows disk;
\item ask for administrator rights, install a driver or a service.
\end{itemize}
\end{securitebox}

\section{Multi-slot card readers}

A reader that takes both an SD and a CF card appears as a single device with several
letters. Windows can only eject the whole reader: Pipit reminds you before doing so.

\capture[0.6]{../capture_lecteur-en.png}

Tick \menu{Don't ask again} if you only ever use one slot at a time.

\section{Settings}\label{sec:reglages}

\menu{Settings…} in the icon's menu opens this window. Closing it does not quit Pipit:
the icon stays next to the clock. To really quit: \menu{Quit Pipit} in the menu.

\capture{../capture_reglages-en.png}

\subsection{Connected devices}
\begin{itemize}
\item \textbf{Displayed name}: double-click to give a meaningful name (“Photo backup”
  rather than “SanDisk Extreme Pro USB Device”). Clear the name to go back to the
  hardware one.
\item \textbf{Hide}: the device no longer appears in the menu (a drive you never unplug,
  for instance). Names and hiding are remembered per hardware model: Pipit remembers
  them even while the device is unplugged.
\end{itemize}

\subsection{Keyboard shortcuts}
\begin{itemize}
\item \textbf{Open the menu}: \touche{Win}+\touche{Alt} by default (the same as in
  USB Safely Remove), or another of your choice.
\item \textbf{Eject all}: none by default, to avoid an accidental ejection.
\end{itemize}
If another program already uses the chosen shortcut, Pipit says so under the list:
pick another one.

\subsection{Options}
\begin{itemize}
\item \textbf{Notify when it is safe to unplug}: the confirmation balloon.
\item \textbf{Ask before ejecting a multi-card reader}: see above.
\item \textbf{Also show hot-plug internal disks}: Windows often presents internal SATA
  disks as “removable”. Pipit hides them so that an internal disk is not ejected by
  mistake; tick only if you have a removable drive bay.
\item \textbf{Start Pipit with Windows}.
\item \textbf{Language}: French, English, or that of Windows.
\end{itemize}

\section{Updates and incident reports}

On start, Pipit quietly checks whether a new version exists, and offers it. Every update
is \textbf{signed}: before installing it, Pipit checks that it really comes from its
author and was not altered on the way.

If Pipit crashes or freezes, and only if you agreed, a technical report is sent
automatically (version, system, error trace; neither your files nor the list of your
devices). \menu{Report a problem…} lets you write one yourself.

\section{Command line}

From a command prompt, in Pipit's folder:

\begin{center}
\begin{tabularx}{\linewidth}{>{\ttfamily\small}l X}
\toprule
\textrm{\textbf{Command}} & \textbf{Effect} \\
\midrule
run.py --list & lists what can be removed, without ejecting anything \\
run.py --eject E & ejects the device holding drive E: \\
run.py --discret & starts the icon only, without opening the settings \\
\bottomrule
\end{tabularx}
\end{center}

\section{Frequently asked questions}

\textbf{My external drive does not show up.} If it is connected through eSATA or a SATA
bay, tick “Also show hot-plug internal disks”. If it is hidden, untick \menu{Hide} in the
settings.

\textbf{Do mice, keyboards or phones show up?} No: Pipit only lists drives (sticks, disks,
memory cards). Everything else can be unplugged without precaution.

\textbf{After ejection, the drive keeps spinning.} Depending on the enclosure, power is not
cut; it is nonetheless properly unmounted. You can unplug it.

\textbf{I want an ejected device back without unplugging it.} Unplug it and plug it back
in: Windows cannot “re-insert” an ejected device.

\end{document}
